\documentclass[../main.tex]{subfiles}

\begin{document}


\chapter{ Week 2 }

代靖涵 25120222201319

\begin{exercise}{}{}
    试作\(  (-1,1)  \)到 \(  \left[ -1,1 \right]   \)之间的一个双射.  
\end{exercise}
\begin{proof}
    取(-1,1)的一个可数 子集 \(  \left\{ \frac{1 }{n }  \right\}_{n\in \mathbb{Z} _{\ge 2}}  \). 做映射 \(  f: \left\{ -1,1 \right\}\cup \left\{ \frac{1 }{n }  \right\}_{n \in \mathbb{Z} _{\ge 2}}\to \left\{ \frac{1 }{n }  \right\}_{n\in \mathbb{Z} _{\ge 2}} \), \[
    f\coloneqq \begin{cases} \frac{1 }{2 },& x= -1\\
    \frac{1 }{3 },&x = 1\\ 
     \frac{1 }{n+ 2 },& x= \frac{1 }{n }    ,n \in \mathbb{Z} _{\ge 2} \end{cases} 
    \]  则 \(  f  \)是一个双射. 定义 \(  F:\left[ -1,1 \right]\to \left( -1,1 \right)    \) \[
    F\left( x \right)\coloneqq \begin{cases} f\left( x \right),& x \in \left\{ -1,1 \right\}\cup \left\{ \frac{1 }{n }  \right\}_{n \in \mathbb{Z} _{\ge 2}}\\ 
     x,& x \in \left( -1,1 \right) \setminus \left\{ \frac{1 }{n }  \right\}_{n\in \mathbb{Z} _{\ge 2}}  \end{cases}  
    \]  则 \[
    \operatorname{Im}F= \operatorname{Im}f\cup \left( \left( -1,1 \right)\setminus \left\{ \frac{1 }{n }  \right\}_{n\in  \mathbb{Z} _{\ge 2}}  \right)= \left\{ \frac{1 }{n }  \right\}_{n\in \mathbb{Z} _{\ge 2}}\cup \left( \left( -1,1 \right)\setminus \left\{ \frac{1 }{n }  \right\}_{n\in \mathbb{Z} _{\ge 2}}  \right)= \left( -1,1 \right)   
    \]故 \(  F  \) 是一个满射. 此外, 任取 \(  x,y \in \left[ -1,1 \right]   \), 设 \(  F\left( x \right)= F\left( y \right)    \), 则 若 \(  F\left( x \right)= F\left( y \right)\in \left\{ \frac{1 }{n }  \right\}_{n \in \mathbb{Z} _{\ge 2}}    \), 则  \(  f\left( x \right)= f\left( y \right)\implies x= y    \); 若 \(  F\left( x \right)= F\left( y \right)\in \left( -1,1 \right)\setminus \left\{ \frac{1 }{n }  \right\}_{ n\in \mathbb{Z} _{\ge 2}}     \), 则 \(  x= F\left( x \right)= F\left( y \right)= y    \). 因此 \(  F  \)是一个单射. 最终我们得到 \(  F  \)是 \(  \left[ -1,1 \right]   \)到 \(  (-1,1)  \)的一个双射.     逆映射 \(  F^{-1}   \)即为题目所需的双射.       
\end{proof}

\begin{exercise}{}{}
    试证明系数为有理数的多项式的全体组成一个可数集.
\end{exercise}
\begin{proof}
   定义 \[
   \mathbb{Q} ^{k}[x]\coloneqq\left\{  \sum _{m= 0}^{k}a_{m}x^{m}: a_0,a_1,\cdots ,a_{m}\in \mathbb{Q}  \right\}
   \]易见存在 \(  \mathbb{Q} ^{k}\left[ x \right]   \)到 \(  \mathbb{Q} ^{k+ 1}  \)之间的一个双射. 注意到 \[
   \mathbb{Q} [x]= \bigcup _{k= 0}^{\infty}\mathbb{Q} ^{k}[x]
   \]  其中每个 \(  \mathbb{Q} ^{k}\left[ x \right]   \)都与 \(  \mathbb{Q} ^{k+ 1}  \)等势. 而 \(  \mathbb{Q} ^{k+ 1}  \)作为可数集的笛卡尔积仍是可数集. 又可数集的可数并仍可数, 因此 \(  \mathbb{Q} \left[ x \right]   \)也是一个可数集.   
\end{proof}
\begin{exercise}{}{}
设 $E \subset (0,1)$ 是无限集. 若从 E 中任意选取不同的数所组成的无穷正项级数总是收敛的, 试证明 E 是可数集.
\end{exercise}
\begin{proof}
令 \[
E_{n}\coloneqq E\cap (\frac{1 }{n+ 1 },\frac{1 }{n }  ]
\]则 \[
E= \bigcup _{n = 1}^{\infty}E_{n}
\] 若存在 \(  k \in \mathbb{Z} _{> 0}  \)使得 \(  E_{k}  \)为不可数集, 则 \(  E_{k}  \) 包含一个可数无限子集, 将其写成一个序列 \(  \left\{ a_{j} \right\}_{j \in \mathbb{N} }  \).   由于 \(  \lim_{j\to \infty}a_{j}= 0  \)不成立, 故级数\[
\sum _{k= 1}^{\infty}a_{k}
\] 发散, 与题设矛盾.  因此, 对于任意的 \(  k \in \mathbb{Z} _{> 0}  \), \(  E_{k}  \)可数, 进而 \(  E  \)是可数集.   
\end{proof}
\begin{exercise}{}{}
试作开圆 $\{(x,y): x^2+y^2<1\}$ 与闭圆盘 $\{(x,y): x^2+y^2 \le 1\}$ 之间的一一对应.
\end{exercise}
\begin{proof}
    记 \[
    D\coloneqq \left\{ \left( x,y \right):x^{2}+ y^{2}< 1  \right\},\quad  \overline{D}\coloneqq \left\{ \left( x,y \right):x^{2}+ y^{2}\le 1  \right\}
    \]令 \[
    S\coloneqq   \overline{D}\setminus D= \left\{ \left( x,y \right):x^{2}+ y^{2}= 1  \right\}
    \] 只需要找到 \(  D\setminus \left\{ \left( 0,0 \right)  \right\}  \)与 \(  \bar{D}\setminus \left\{ \left( 0,0 \right)  \right\}  \)之间的一一对应.  将圆盘上的点用极坐标表示, 我们寻找 \[
    (0,1)\times \left[ 0,2 \pi \right)\to  (0,1]\times \left[ 0,2 \pi \right)  
    \]的一一对应. 令 \(  f:(0,1]\to (0,1)  \) \[
    f\left( x \right)= \begin{cases} \frac{1 }{n+ 1 },&x= \frac{1 }{n }, n\in \mathbb{Z} _{> 0}  \\ 
     x, & x \in (0,1]\setminus \left\{ \frac{1 }{n }  \right\}_{n\in \mathbb{Z} _{> 0}} \end{cases}  
    \] 则 \(  f  \)是一个双射. 定义 \(  F:(0,1]\to [0,2 \pi)\to (0,1)\times [0,2 \pi]  \) \[
    F\left( r, \theta  \right)= \left( f\left( r \right), \theta   \right)  
    \]则 \(  F  \)是一个双射.   令 \(  \rho \left( r, \theta  \right)= \left( r\cos  \theta , r\sin  \theta  \right)    \)为极坐标映射, 它是 \(  \mathbb{R} ^{2}\setminus \left\{ 0,0 \right\}  \)上的同胚映射. 最后, 令 \(  G:\bar{D}\to D \) \[
    G\left( x,y \right)= \begin{cases}\left( \rho \circ  F\circ  \rho ^{-1}   \right)\left( x,y \right),& \left( x,y \right)\neq \left( 0,0 \right)\\ 
     \left( 0,0 \right),& \left( x,y \right)= \left( 0,0 \right)        \end{cases}  
    \] 则 \(  G  \)是一个双射. 
\end{proof}
\begin{exercise}{}{}
    设 $f(x)$ 是定义在 $[0,1]$ 上的实值函数, 且存在常数 $M$, 使得对于 $[0,1]$ 中任意有限个数: $x_1, x_2, \cdots, x_n$, 均有
\[ |f(x_1) + f(x_2) + \cdots + f(x_n)| \le M, \]
试证明下述集合是可数集: $E = \{x \in [0,1]: f(x) \ne 0\}$.
\end{exercise}

\begin{proof}
令 \[
E^{+ }= f^{-1} \left( 0,\infty \right) ,\quad E^{-}= f^{-1} \left( (-\infty,0) \right) 
\]则 \[
E= E^{+ }\cup E^{-}
\]注意到\[
  E^{+ }= f^{-1} \left( (0,\infty) \right)= \bigcup _{n = 1}^{\infty}f^{-1} \left( \left( \frac{1 }{n },\infty  \right) \right)  
  \]令 \(  E_{n}= f^{-1} \left( \left( \frac{1 }{n },\infty  \right) \right)   \). 若存在 \(  n_0\in \mathbb{Z} _{> 0}  \), 使得 \(  E_{n_0}  \)是不可数集. 对于题目给定的 \(  M  \), 找到使得  \[
  k_{M}\frac{1 }{n_0 }> M 
  \]的正整数 \(  k_{M}  \)      . \(  E_{n_0}  \)存在一个含有 \(  k_{M}  \)个元素的子集 \(  x_1,x_2,\cdots ,x_{k_{M}}  \). 满足 \[
  \left| f\left( x_1 \right)+ \cdots + f\left( x_{k_{M}} \right)   \right|>  k_{M}\frac{1 }{n_0 }> M  
  \], 与题设矛盾. 因此对于任意的 \(  n\in \mathbb{Z} _{> 0}  \), \(  E_{n}  \)都是一个可数集. 进而 \(  E^{+ }  \)是一个可数集.      类似地, 可以证明 \(  E^{-}  \)是一个可数集, 进而 \(  E= E^{+ }\cup E^{-}  \)是一个可数集.  
\end{proof}
 \begin{exercise}{}{}
设 $f(x)$ 是定义在 $\mathbf{R}^1$ 上的实值函数. 如果对于任意的 $x_0 \in \mathbf{R}^1$, 必存在 $\delta > 0$, 当 $|x - x_0| < \delta$ 时, 有 $f(x) \ge f(x_0)$, 试证明集合 $E = \{y: y = f(x)\}$ 是可数集.
\end{exercise}
\begin{proof}
    令 \[
    A_{n}\coloneqq \left\{ x_0\in \mathbb{R} : \text{当} \left| x-x_0 \right|< \frac{1 }{n }\text{时, 有} f\left( x \right)\ge f\left( x_0 \right)     \right\}
    \] 则 由题设可知 \[
    \mathbb{R} = \bigcup _{n = 1}^{\infty}A_{n}
    \]再令 \[
    A_{n,k}=  A_{n}\cap (\frac{k }{n },\frac{k+ 1 }{n }  ]
    \]则 \[
   A_{n}= \bigcup _{k= -\infty}^{\infty}A_{n,k},\quad  \mathbb{R} = \bigcup _{n = 1}^{\infty}\bigcup _{k = -\infty}^{\infty}A_{n,k}
    \]令 \[
    E_{n,k}\coloneqq  f\left( A_{n,k} \right) 
    \]则 \[
    E= \bigcup _{n = 1}^{\infty}\bigcup _{k= -\infty}^{\infty}E_{n,k}
    \] 若 \(  E_{n,k}  \)非空, 则取 \(  y_1\in E_{n,k}  \). 存在 \(  x_1 \in A_{n}\cap (\frac{k }{n },\frac{k+ 1 }{n }  ]  \), 使得 \(  f\left( x_1 \right)= y_1   \). 任取 \(  y_2 \in E_{n,k}  \), 存在 \(  x_2\in A_{n}\cap (\frac{k }{n },\frac{k+ 1 }{n }  ]  \), 使得 \(  f\left( x_2 \right)= y_2   \)       . 注意到 \(  \left| x_1-x_2 \right|< \frac{1 }{n }    \), 根据 \(  A_{n}  \)的定义, 我们有 \(  f\left( x_1 \right)\ge f\left( x_2 \right)    \)以及 \(  f\left( x_2 \right)\ge f\left( x_1 \right)    \). 因此 \(  y_1= y_2  \), 此时 \(  E_{n,k}  \)为单点集 \(  \left\{ y_1 \right\}  \).        于是对于所有的 \(  n \in \mathbb{Z} _{> 0}  \),  \[
    \bigcup _{k= -\infty}^{\infty}E_{n,k}
    \]是单点集或空集的可数并, 从而是至多可数集.  进而 \(  E  \)作为可数集的可数并也是一个可数集. 
\end{proof}

\begin{exercise}{}{}
    设 $f(x)$ 在 $(a,b)$ 上可微, 且除可数集外, 有 $f'(x)=0$, 试证明 $f(x)=c$(常数).
\end{exercise}

\begin{proof}
    若 \(  f  \)不为常值, 则 \(  f^{\prime}   \)在 \(  \left( a,b \right)   \)上不恒为零.
    \begin{enumerate}
        \item 若 \(  f^{\prime}   \)在 \(  (a,b)  \)上取非零常值, 此时 \(  \left\{ f^{\prime} \left( x \right)\neq 0  \right\}= \left( a,b \right)   \)为不可数集.
        \item 若 存在 \(  x_1  \), \(  x_2  \)使得 \(  f^{\prime} \left( x_1 \right)\neq f^{\prime} \left( x_2 \right)    \), 不妨设 \(  x_2> x_1  \), \(  f^{\prime} \left( x_2 \right)> f^{\prime} \left( x_1 \right)    \) .       此时 \(  f  \)为 \(  \left[ x_1,x_2 \right]   \)上的可微函数,   由 Darboux定理, 对于任意的 \(  c \in \left( f^{\prime} \left( x_1 \right),f^{\prime} \left( x_2 \right)   \right)   \),  都 存在 \(  x_{c}  \)使得 \(  f^{\prime} \left( x_{c} \right)= c  \). 特别地, 我们可以将 \(  c  \)的取值限制在 \(  \left( f^{\prime} \left( x_1 \right),f^{\prime} \left( x_2 \right)   \right)\setminus \left\{ 0 \right\}   \)上. 即存在单射 \(  F: \left( f^{\prime} \left( x_1 \right),f^{\prime} \left( x_2 \right)   \right)\setminus \left\{ 0 \right\}\to \left\{ f^{\prime} \left( x \right)\neq 0  \right\}   \) \[
        F\left( c \right)= x_{c} 
        \]   这表明 \[
        \left| \left\{ f^{\prime} \left( x \right)\neq 0  \right\} \right|\ge \left| \left( f^{\prime} \left( x_1 \right),f^{\prime} \left( x_2 \right)   \right)\setminus \left\{ 0 \right\}  \right|  
        \]其中 \(  \left( f^{\prime} \left( x_1 \right),f^{\prime} \left( x_2 \right)   \right)   \)有连续统的势, 为不可数集, 从而 \(  \left( f^{\prime} \left( x_1 \right),f^{\prime} \left( x_2 \right)   \right) \setminus \left\{ 0 \right\}  \)为不可数集, 进而 \(  \left\{ f^{\prime} \left( x \right)\neq 0  \right\}  \)也不可数.   
    \end{enumerate}
       至此, 我们说明了如果 \(  f\left( x \right)   \)不为常值函数, 则 \(  \left\{ f^{\prime} \left( x \right)\neq 0  \right\}  \)一定不可数, 因此原命题成立.  
\end{proof}

\begin{exercise}{}{9-11-1}
    设 $E \subset \mathbf{R}$ 是不可数集, 则 $E' \neq \varnothing$.
\end{exercise}
\begin{proof}
    若 \(  E  \)为不可数集, 则存在 \(   n\in \mathbb{Z} _{> 0}  \), \(  \left[ -n,n  \right]\cap E   \)为不可数集. 若不然, \[
    E = \bigcup  _{n = 1}^{\infty} \left( \left[ -n,n \right]\cap E  \right) 
    \]   为可数集的可数并, 进而 \(  E  \)是至多可数的,  矛盾. 下设 \(  n_0 \in \mathbb{Z} _{> 0}  \), \(  E_{n_0}\coloneqq \left[ -n_0,n_0 \right]\cap E   \)   为不可数集. 即 \(  I_0= [-n_0, n_0]  \). 将 \(  I_0  \)二等分为 \(  I_0= I_1\cup I_2\coloneqq [-n_0,0]\cup [0,n_0]  \).则 记 \(  E_{n_0,i}\coloneqq E_{n_0}\cap I_{i}, i= 1,2  \). 则 \(  E_{n_0,1}, E_{n_0,2}  \)中至少一个为不可数集, 取其一记作 \(  E_{n_0,i_1}  \).      闭区间 \(  I_{i_1}  \) 二分为 \(  I_{i_1,1}, I_{i_1,2}  \), 记 \(  E_{n_0,i_1,j}\coloneqq I_{i_1,j}\cap E, j= 1,2  \). 类似地,  \(  E_{n_0,i_1,1}, E_{n_0,i_1,2}  \)中至少一个为不可数集, 取其一记作 \(  E_{n_0,i_1,i_2}  \). 反复上述过程, 若我们得到不可数集 \(  E_{n_0,i_1,\cdots ,i_{n}}  \), 可以将区间 \(  I_{i_{n}}  \)二分, 类似地得到不可数集 \(  E_{n_0,i_1,\cdots ,i_{n+ 1}}  \).       此时, 我们有闭区间套 \[
    I_{i_1}\supseteq I_{i_1,i_2}\supseteq \cdots I_{i_1,i_2,\cdots ,i_{n}}\supseteq  \cdots 
    \]由闭区间套定理, 存在唯一的点 \(  x    \), 使得 \[
     x \in I_{i_1,\cdots i_{n}},\quad \forall  n\in \mathbb{N} 
    \] 首先, 取定 \(  x_0 \in E_{n_0}  \). 对于任意的 \(  k \in \mathbb{Z} _{> 0}  \) , 记 \(  B_{k}= \left( x-\frac{1 }{k },x+ \frac{1 }{k }   \right)   \)  . 若现有两两不同的 \(  x_0,x_1,\cdots ,x_{k-1}  \), 使得 \(  x_{j}\in B_{j}\cap E, j= 0,1,\cdots ,k-1  \).   考虑到 存在 \(  n_{k}  \)使得 \(  B_{k}\supseteq I_{i_1,\cdots ,i_{n_{k}}}  \supseteq E_{n_0,i_1,\cdots ,i_{n_{k}}}\)  . 而 \(  E_{n_0,i_1,\cdots ,i_{n_{k}}}  \) 为无限点集,  故存在其中的异于 \(  x_0,x_1,\cdots ,x_{k-1}  \)的点 \(  x_{k}  \), 且使得 \(  x_{k}\in B_{k}\cap E  \)   . 我们递归地构造了一个两两不同点列 \(  \left\{ x_{n} \right\}  \), 使得 \[
    \left\{ x_{n} \right\}\subseteq E,\quad \lim _{n\to \infty}x_{k}= x
    \]这表明 \(  x \in E^{\prime}   \), 因此 \(  E^{\prime} \neq \varnothing  \).   
\end{proof}

\begin{exercise}{}{}
设 $F \subset \mathbb{R}$ 是有界集, $f(x)$ 是定义在 $F$ 上的(实值)函数. 若对任意的 $x_0 \in F'$, 均有 $f(x) \to +\infty (x \in F \text{ 且 } x \to x_0)$, 试证明 $F$ 是可数集. ($F = \bigcup_{n=1}^{\infty} \{x: f(x) \le n\}$)
\end{exercise}



\begin{proof}
    由于 \[
    F=  \bigcup _{n = 1}^{\infty}\left\{ x:f\left( x \right)\le n  \right\}
    \]记 \(  \left\{ x:f\left( x \right)\le n  \right\}\coloneqq  F_{n}  \), 只需证明对于每个 \(  n  \),  \(  F _{n} \)是可数集.
    现任取定正整数 \(  n  \).   若 \(  F_{n}^{\prime} \neq \varnothing  \), 则可 取 \(  x \in F_{n}^{\prime}   \subseteq F^{\prime} \), 存在 \(  F_{n}  \) 上两两互异的点列 \(  \left\{ x_{n} \right\}  \), 使得 \[
    \lim_{n\to \infty}x_{n}= x
    \] 由题设可知 \[
    \lim_{n\to \infty}f\left( x_{n} \right)= \infty 
    \]从而存在 \(  N  \), 使得 \(  f\left( x_{N} \right)>  n  \), 这与 \(  F_{n}  \)的定义矛盾.    因此 \(  F^{\prime}_{n} = \varnothing \).  此时若 \(  F_{n}  \)不可数, 将与Exercise \ref{ex:9-11-1}矛盾. 又 \(  n  \)是任取的,  \(  F  \)作为可数集的可数并也是可数的.
    
    \begin{remark}
        这里借助Exercise \ref{ex:9-11-1}完成了一个不需要 \(  F  \)有界的 该命题的证明.
        
        如果不借助Exercise \ref{ex:9-11-1} 并假设 \(  F  \)有界. 此时\(  F_{n}  \)作为 \(  F  \)的子集也是有界的.   可以通过 \(  \mathbb{R}   \)上的Bolzano-Weierstrass定理, 利用 事实``有界无限点集存在极限点''导出更强的结果: \(  F_{n}  \)是一个有限集.
    \end{remark}
\end{proof}

\begin{exercise}{}{}
设 $\{f_\alpha(x)\}_{\alpha \in I}$ 是定义在 $[a,b]$ 上的实值函数族. 若存在 $M > 0$, 使得
\[\begin{aligned}
|f_\alpha(x)| \le M, \quad x \in [a,b], \quad \alpha \in I,
\end{aligned}\]
试证明对 $[a,b]$ 中任一可数集 $E$, 总有函数列 $\{f_{\alpha_n}(x)\}$, 存在极限
\[\begin{aligned}
\lim_{n \to \infty} \{f_{\alpha_n}(x)\}, \quad x \in E.
\end{aligned}\]
\end{exercise}
\begin{remark}
    或许需要指定 \(  I  \)是无限集, 且要求序列 \(  \left\{ \alpha _{n} \right\}  \) 包含无限多个元素/不允许重复? 否则命题或是平凡的, 或是不成立的.  下面的证明中默认了这种设置.
\end{remark}
\begin{proof}
    记 \(  E= \left\{ x_1,x_2,\cdots  \right\}  \) . 对于 \(  x_1  \), 考虑有界的实数集 \(  \left\{ f_{\alpha }\left( x_1 \right)  \right\}_{\alpha \in I}  \), 可以取出可数的指标序列 \(  \left\{ \alpha _{1,n} \right\}_{n = 1}^{\infty}  \), 使得序列 \(  \left\{ f_{ \alpha _{1,n}}\left( x_1 \right)  \right\}_{n = 1}^{\infty}  \)收敛.  记 \(  I_1  \)为指标序列的值集 . 考虑有界的实数族 \(  \left\{ f_{\alpha }\left( x_2 \right)  \right\}_{\alpha \in I_1}  \). 同样地, 可以取出\(  I_1  \)中 可数的指标序列 \(  \left\{ \alpha _{2,n} \right\}_{n = 1}^{\infty}  \), 使得 它是 \(  \left\{ \alpha _{1,n} \right\}_{n = 1}^{\infty}  \)的子列, 且  序列 \(  \left\{ f_{\alpha _{2,n}}\left( x_2 \right)  \right\}_{n = 1}^{\infty}  \)收敛. 以此类推, 若我们得到收敛的序列 \(  \left\{ f_{\alpha _{j,n}}\left( x_{j} \right)  \right\}_{n = 1}^{\infty}\),以及指标集 \(  I_{j}  \), \(  j= 1,2,\cdots ,k -1 \).  可以从有界实数族 \(  \left\{ f_{\alpha }\left( x_{k} \right)  \right\}_{ \alpha \in I_{k-1}}  \)    中取出可数的指标序列 \(  \left\{ \alpha _{k,n} \right\}_{n = 1}^{\infty}  \) ,它是 \(  \left\{ \alpha _{k-1,n} \right\}_{n = 1}^{\infty}  \)的一个子列,  使得序列 \(  \left\{ f_{\alpha _{k,n}}\left( x_{k} \right)  \right\}_{n = 1}^{\infty}  \)收敛. 现在, 重新定义新的指标序列 \[
    \left\{ a_{n} \right\}, \quad a_{n}\coloneqq \alpha _{n,n}
    \]则对于任意的 \(  x_{k } \in E  \), 极限 \[
    \lim_{n\to \infty}f_{a_{n}}\left( x_{k} \right)= \lim_{n\to \infty}f_{\alpha _{n,n}} \left( x_{k} \right)
    \]存在. 这是因为序列 \(  \left\{ f_{\alpha _{k,n}}\left( x _{k}\right)  \right\}_{n = 1}^{\infty}  \)  收敛, 且\(  \left\{ \alpha  _{n,n} \right\}_{n = k}^{\infty}  \)是 \(  \left\{ \alpha _{k,n} \right\}_{n = 1}^{\infty}  \)的子列, 而收敛序列的任意子列也是收敛的.   
\end{proof}
\end{document}